\documentclass[12pt,letterpaper]{article}
\usepackage[UTF8]{ctex}
\usepackage{fullpage}
\usepackage[top=2cm, bottom=4.5cm, left=2.5cm, right=2.5cm]{geometry}
\usepackage{amsmath,amsthm,amsfonts,amssymb,amscd}
\usepackage{lastpage}
\usepackage{enumerate}
\usepackage{fancyhdr}
\usepackage{mathrsfs}
\usepackage{xcolor}
\usepackage{graphicx}
\usepackage{listings}
\usepackage{hyperref}
\usepackage{dot2texi}
\usepackage{tikz}
\usepackage{float}
\usepackage{pifont}
\usepackage[pdf]{graphviz}
\usetikzlibrary{automata,shapes,arrows} 

\hypersetup{%
  colorlinks=true,
  linkcolor=blue,
  linkbordercolor={0 0 1}
}
 
\renewcommand\lstlistingname{Algorithm}
\renewcommand\lstlistlistingname{Algorithms}
\def\lstlistingautorefname{Alg.}
\renewcommand\refname{参考文献} 

\lstdefinestyle{Matlab}{
    language        = matlab,
    frame           = lines, 
    basicstyle      = \footnotesize,
    keywordstyle    = \color{blue},
    stringstyle     = \color{green},
    commentstyle    = \color{red}\ttfamily
}
\setlength{\parindent}{0.0in}
\setlength{\parskip}{0.05in}

\pagestyle{fancyplain}
\headheight 35pt
\lhead{\name\\\ID}                 
\chead{\textbf{\Large \hwtopic}}
\rhead{\course \\ \today}
\lfoot{}
\cfoot{}
\rfoot{\small\thepage}
\headsep 1.5em



% Edit these as appropriate
\newcommand\course{FNLP 2024}
\newcommand\hwtopic{Assignment 2}
\newcommand\name{张三}                % <-- Name
\newcommand\ID{220xxxxxxx}           % <-- Student ID

\begin{document}

\begin{center}
    \LARGE
    \bf
    POS Tagging
\end{center}

\section{引言}
报告可用中文或者英文撰写。

在需要引用文献时，请通过Google Scholar、Semantic Scholar、ACL Anthology等网站获取文献的BibTeX信息，放入\texttt{ref.bib}文件中，然后通过\texttt{cite}命令引用该文献~\cite{vaswani2017attention}.

\section{XXX}


\section{结论}


\section{AI使用声明}
\textbf{报告必须包含此节，并准确填写相关内容，否则将会扣分。}

本次作业的目标是独立理解并实现一个POS Tagging系统。允许在有限、辅助性的范围内使用 AI 工具，但不得替代你自身的实现与思考过程。注意：将LLM作为方法或模型的一部分（例如设计基于 LLM 的 POS Tagging 方法）属于算法设计选择，不属于本节讨论的AI使用范畴。

允许的AI使用（需如实披露）：
\begin{itemize}
    \item 询问概念性问题（如让AI解释 POS tagging、HMM、Viterbi 算法等）
    \item 寻找相关文献并阅读（AI可能因幻觉给出不存在的文献，请注意甄别）
    \item 在写代码时使用Copilot等工具进行基本的代码提示
    \item 对少量代码片段进行debug
    \item 改进报告的语言表达
\end{itemize}


禁止的AI使用：
\begin{itemize}
    \item 上传或提供作业说明（如完整的PDF或文本形式的作业要求）给 AI，并要求其生成完整或接近完整的解决方案
    \item 使用 AI 生成代码的核心实现部分（如解码算法、特征设计等）
    \item 使用 AI 完成报告的内容撰写
    \item 未经自己理解和验证的情况下直接使用AI生成的内容
\end{itemize}

请阅读以上须知后，填写以下内容：

\subsection{AI 使用情况（单选）}
（对于你选择的那项，请将对应那行的latex代码中的$\backslash$square替换成$\backslash$blacksquare） 

\begin{itemize}
    \item[$\square$] 未使用任何 AI 工具
    % \item[$\blacksquare$] 未使用任何 AI 工具
    \item[$\square$] 仅在允许范围内使用 AI 工具
    % \item[$\blacksquare$] 仅在允许范围内使用 AI 工具
    \item[$\square$] 使用了超出允许范围的 AI 工具（需详细说明，可能会影响评分或按学术不端处理）
    % \item[$\blacksquare$] 使用了超出允许范围的 AI 工具（需详细说明，可能会影响评分或按学术不端处理）
\end{itemize}

\subsection{使用情况说明}
如使用过 AI，请说明：
\begin{itemize}
    \item 你使用了何种AI工具，涉及作业的哪些部分
    \item 你如何对 AI 生成内容进行验证与修改
\end{itemize}


\subsection{承诺确认}
（请将latex代码中“本人确认”那行的$\backslash$square替换成$\backslash$blacksquare）
\begin{itemize}
    \item[$\square$] 本人确认：
    % \item[$\blacksquare$] 本人确认：
    \begin{itemize}
        \item 已如实披露所有 AI 使用情况
        \item 未以任何禁止方式使用 AI
        \item 已理解并能够解释提交作业的全部内容，在助教抽查的情况下，可以解释提交作业中的任何部分
    \end{itemize}
\end{itemize}


\bibliographystyle{plain}
\bibliography{ref}

\end{document}
